\documentclass[11pt]{article}
\usepackage{amsmath,amssymb,amsthm,mathrsfs,hyperref,geometry,booktabs}
\geometry{margin=1in}
\newtheorem{theorem}{Theorem}[section]
\newtheorem{lemma}[theorem]{Lemma}
\newtheorem{proposition}[theorem]{Proposition}
\newtheorem{corollary}[theorem]{Corollary}
\theoremstyle{definition}
\newtheorem{definition}[theorem]{Definition}
\newtheorem{hypothesis}{Cited Proposition}
\theoremstyle{remark}
\newtheorem{remark}[theorem]{Remark}
\newcommand{\T}{\mathfrak{T}_F}
\newcommand{\ip}[2]{\langle #1,#2\rangle}
\newcommand{\lean}[1]{\texttt{\detokenize{#1}}}
\newcommand{\LeanRef}[1]{\par\noindent\emph{Lean:} \lean{#1}.\par}
\title{The Frohmanian Symplectic Tether $\T$:\\ Written Proofs Companion to the Lean~4 Formalization}
\author{Benjamin Stanley Frohman\\ \small ORCID \href{https://orcid.org/0009-0006-7068-3718}{0009-0006-7068-3718}}
\date{October 29, 2026}
\begin{document}
\maketitle
\begin{abstract}
This is a sibling paper to the Lean~4 repository \texttt{frohmanian-symplectic-tether-may-form}.
Each result is given a written proof and cross-referenced to the Lean declaration that
kernel-checks it. Result~1 is geometric: under three admissibility conditions, the symmetric
correction $\T[\kappa,w,u](F,G)=-(\kappa w)\ip{\Pi_uF}{\Pi_uG}$ is unique up to scale, and
Route~A fixes the scale. Result~2 is about a \emph{modified} system, the tethered vorticity equation
$\partial_t\omega+(u\cdot\nabla)\omega=(\omega\cdot\nabla)u+\nu\Delta\omega-\kappa P(|\omega|^2\omega)$.
For that system we prove the $L^2$ bound, the $H^1$ estimate, the log-loss comparison and global
smoothness, \emph{conditional on cited propositions that are listed explicitly}.
\textbf{None of this resolves the Clay Millennium problem.} Plain Navier--Stokes is unchanged
(Section~\ref{sec:open}).
\end{abstract}

\section{Physical motivation}\label{sec:motivation}
The starting picture is physical. The vortex-stretching term $(\omega\cdot\nabla)u$ behaves as if
a hidden geometric constraint, an ``invisible object'', were shaping three-dimensional vorticity. The
programme has two steps. (i) \emph{Derive} what such an object can be from symmetry and
energy considerations alone, in Hamiltonian/Lie--Poisson language. (ii) \emph{Insert} the derived
object into the dynamics, at vorticity level, through the Leray projection $P$.

Two pictures from the author's June~22, 2026 compilation guide the derivation. Both are formalized:
\begin{itemize}
\item \textbf{Chevalley--Eilenberg 2-cocycle.} Antisymmetric corrections to a Lie--Poisson bracket are
CE 2-cochains. Jacobi requires them to be cocycles, and $d\circ d=0$ makes coboundaries automatic
cocycles. A \emph{symmetric} correction is not a CE cochain at all. That is why the tether is
metriplectic and not Hamiltonian (Prop.~\ref{prop:ce}).
\item \textbf{$T^3$ shear.} The shear $(\sin y,0,0)$ is divergence free, so it is admissible. The
variant $(\sin x,0,0)$ is not (its divergence is $1$ at the origin). This is the
counterexample that shows the divergence-free constraint is essential (Prop.~\ref{prop:shear}); at field level it is essential but not sufficient for Jacobi (Theorem~\ref{thm:fieldjacobi}).
\end{itemize}
Whether real fluids carry the resulting term is \emph{not} decided here (Section~\ref{sec:open}).

\section{Result 1: geometric uniqueness of $\T$}\label{sec:geom}
Throughout, $H$ is a real inner-product space with $\dim H\ge 3$, $u\in H$, $u\neq0$, and
$\Pi_u x = x-\frac{\ip{u}{x}}{\|u\|^2}u$.
\begin{definition}
$\T[\kappa,w,u](F,G):=-(\kappa w)\ip{\Pi_uF}{\Pi_uG}$. \LeanRef{Frohman.Tether.frohmanianTether; frohmanianTether_apply}
\end{definition}

\begin{lemma}[Basic properties]
(a) $\T$ is symmetric. (b) $\T(F,F)\le0$ whenever $\kappa w\ge0$. (c) It is degenerate with respect to the
kinetic-energy gradient: $\T(F,u)=0$ for all $F$. (d) It is invariant under every isometry
commuting with $\Pi_u$.
\end{lemma}
\begin{proof}
(a) is symmetry of the inner product. (b) gives $-\kappa w\|\Pi_uF\|^2$. (c) holds because $\Pi_uu=0$.
(d) If $A$ is an isometry with $A\Pi_u=\Pi_uA$, then $\ip{\Pi_uAF}{\Pi_uAG}=\ip{A\Pi_uF}{A\Pi_uG}=\ip{\Pi_uF}{\Pi_uG}$.
\end{proof}
\LeanRef{frohmanianTether_symm, frohmanianTether_self_nonpos, degeneracy_verification, frohmanianTether_invariant, piU_comm_isometry}

\begin{theorem}[Uniqueness up to scale]\label{thm:unique}
Let $B$ be a symmetric bilinear form that is (C1) invariant under the coadjoint isometries fixing
$u$ and (C2) degenerate with respect to $u$. Then $B=\T[c,1,u]$ for some $c\in\mathbb R$.
If, in addition, (C3) fixes the coefficient ($C3(\T[c,1,u])\Rightarrow c=\kappa w$), then $B=\T[\kappa,w,u]$.
\end{theorem}
\begin{proof}
(C2) and symmetry give $B(x,y)=B(\Pi_ux,\Pi_uy)$ (\lean{degeneracy_forces_projection}).
So $B$ is determined by its restriction to $u^\perp$, which has dimension $\ge2$. Fix a unit $e_0\perp u$
(\lean{exists_unit_orthogonal}). For every unit $e\perp u$, (C1) supplies an isometry fixing $u$
that sends $e_0\mapsto e$, so $B(e,e)=B(e_0,e_0)$ (\lean{diag_eq_of_norm_eq}). By homogeneity,
$B(x,x)=B(e_0,e_0)\|\Pi_ux\|^2$. Polarization recovers $B(x,y)=B(e_0,e_0)\ip{\Pi_ux}{\Pi_uy}$,
that is, $B=\T[c,1,u]$ with $c=-B(e_0,e_0)$. The calibrated statement then follows from (C3).
\end{proof}
\LeanRef{admissible_eq_smul_frohmanianTether, uniqueness_of_minimal_metric_tether, uniqueness_of_minimal_metric_tether_iff}

\begin{theorem}[Route A determines $\kappa$]\label{thm:routeA}
Let $\omega\perp u$, $\|\omega\|=1$, $w\neq0$. Suppose $\T[\kappa,w,u]$ matches the stretching
coefficient $S$ in the sense $\T[\kappa,w,u](\omega,\omega)+S w=0$. Then $\kappa=S$.
Consequently, in $\dim\ge3$ with $S\ge0$, $w>0$, a correction is admissible for Route~A iff
$B=\T[\kappa,w,u]$ with $\kappa=S$.
\end{theorem}
\begin{proof}
Since $\omega\perp u$, $\Pi_u\omega=\omega$, so $\T(\omega,\omega)=-\kappa w\|\omega\|^2=-\kappa w$. The matching gives
$(S-\kappa)w=0$, and $w\ne0$ yields $\kappa=S$. The equivalence combines this with Theorem~\ref{thm:unique}.
\end{proof}
\LeanRef{routeA_determines_kappa, frohmanian_symplectic_tether_theorem, uniqueness_routeA_iff, frohmanianTether_routeA}

\begin{theorem}[Nine-term Jacobi]\label{thm:jacobi}
For a classical (Jacobi) bracket $P$ and the tethered bracket $P+\T$, the Jacobiator expands into
nine terms. The three purely classical ones cancel by Jacobi of $P$, and the three purely tether
terms vanish. So the Jacobiator equals the cyclic sum of the mixed terms, and Jacobi holds on a set
$R$ iff that cyclic sum vanishes on $R$. It vanishes when $u$ is central for $P$, $P$ is
homogeneous in its first slot, and $\Pi_u P(a,b)=0$.
\end{theorem}
\begin{proof}
Bilinear expansion gives the nine terms (\lean{jacobiator_nine_terms}). Terms 7--9 are
$\T$-in-$\T$ and vanish because the tether vector field is $\Pi_u$-valued and annihilated in the
second application (\lean{tetherVec_tetherVec}, \lean{jacobi_term7}--\lean{jacobi_term9}).
The classical triple cancels (\lean{jacobi_classical_cancel}). Under the sufficient hypotheses each
mixed term contains $\Pi_uP(\cdot,\cdot)=0$ or $P(u,\cdot)=0$.
\end{proof}
\LeanRef{jacobiator_eq_mixed, jacobi_iff_mixed_cancel, mixedTerm_eq_zero, jacobiator_eq_zero_of_central}

\begin{theorem}[Metriplectic split]
Every bilinear form splits uniquely into symmetric and antisymmetric parts. $\T$ is purely symmetric,
and the metriplectic bracket $P+\T$ with $P$ antisymmetric conserves kinetic energy:
$\{E,E\}=0$ and $\T(F,\nabla E)=0$.
\end{theorem}
\begin{proof}
$B=\tfrac12(B+B^\top)+\tfrac12(B-B^\top)$. Uniqueness holds because a form that is both symmetric
and antisymmetric is zero. Energy conservation follows from $P(F,F)=0$ and degeneracy.
\end{proof}
\LeanRef{split_unique, symPart_frohmanianTether, antiPart_frohmanianTether, metriplecticBracket_split, metriplectic_energy_conserved}

\begin{proposition}[Chevalley--Eilenberg]\label{prop:ce}
(a) $d_2\circ d_1=0$ on 1-cochains of any Lie algebra. (b) Coboundaries are cocycles.
(c) If $2$ is invertible and the symmetric part of $\beta$ is alternating, it is zero. So a nonzero
symmetric correction such as $\T$ is never a CE 2-cochain.
\end{proposition}
\begin{proof}
(a) $d_2(d_1f)(x,y,z)$ is $-f$ applied to the Jacobi sum $[x,[y,z]]+[y,[z,x]]+[z,[x,y]]=0$.
(b) follows from (a). (c) $s(x,y)=s(y,x)$ and $s(x,y)=-s(y,x)$ give $2s=0$.
\end{proof}
\LeanRef{d2_d1, cocycle_of_coboundary, antiCorr_cocycle_iff, symPart_not_cochain}

\begin{proposition}[$T^3$ shear]\label{prop:shear}
$\mathrm{div}(\sin y,0,0)\equiv0$, whereas $\mathrm{div}(\sin x,0,0)=\cos x$, which equals $1$ at the origin.
\end{proposition}
\LeanRef{shear_divFree, shearBad_div, shearBad_not_divFree}

\subsection{Field-level Jacobi test for the two shears}\label{sec:fieldjacobi}
The algebra above is finite-dimensional, so divergence-freeness never enters it. To let it enter,
we work on trigonometric-polynomial vector fields on $T^3$, written as finite Fourier sums
$a=\sum_k\hat a_k e^{ik\cdot x}$ (\lean{Fld}). On these fields we use:
\begin{itemize}
\item the Lie bracket $[a,b]=(a\cdot\nabla)b-(b\cdot\nabla)a$, mode by mode
$[v e^{ip\cdot x},w e^{iq\cdot x}]=i\bigl((v\cdot q)w-(w\cdot p)v\bigr)e^{i(p+q)\cdot x}$ (\lean{br});
\item the $L^2$ pairing $\langle a,b\rangle=(2\pi)^{-3}\int a\cdot b$ (\lean{pr});
\item the projection $\Pi_ux=x-\frac{\langle u,x\rangle}{\langle u,u\rangle}u$ (\lean{piF});
\item the tether $T(a,b)=-(\kappa w)\langle\Pi_ua,\Pi_ub\rangle\,u$ (\lean{tetF});
\item the tethered bracket $[a,b]+T(a,b)$ (\lean{brT});
\item divergence-free, meaning $k\cdot\hat a_k=0$ for every mode $k$ (\lean{DivFree}).
\end{itemize}
Scalars are complex. The real fields used below have conjugate-symmetric coefficients.
\begin{theorem}[Field-level Jacobi]\label{thm:fieldjacobi}
\begin{enumerate}
\item[(a)] The field bracket satisfies Jacobi. So the Arnold (Lie--Poisson) bracket on linear
functionals is classical.
\item[(b)] If $\langle u,u\rangle\neq0$, the Jacobiator of the tethered bracket equals the cyclic sum
of the mixed terms $[T(a,b),c]+T([a,b],c)$.
\item[(c)] Divergence-free fields along $e_x$ (parallel shears $f(y,z)e_x$) commute pairwise. Hence at
$u=(\sin y,0,0)$ the mixed terms cancel, and Jacobi holds, for every triple in that class.
\item[(d)] At $u=(\sin x,0,0)$, which is not divergence-free, take the divergence-free triple
$a=b=c=(1,0,0)$. The $e^{ix}$, $x$-component coefficient of the mixed sum is $\tfrac32\kappa w$. So
Jacobi fails whenever $\kappa w\ne0$.
\item[(e)] At $u=(\sin y,0,0)$ with the divergence-free triple $a=b=c=(0,0,\sin x)$, the
$e^{i(x+y)}$, $z$-component coefficient is $\tfrac{3i}{8}\kappa w\neq0$. So Jacobi fails here too.
\end{enumerate}
\end{theorem}
\LeanRef{jacF_eq_zero, jacT_eq_mixedSum, br_eq_zero_of_xOnly_divFree, shearY_mixed_cancel, shearY_jacobi_holds, not_divFree_uX, shearX_mixed_coeff, shearX_jacobi_fails, shearY_transverse_coeff, shearY_jacobi_fails_transverse}
\noindent\emph{Scope.} Divergence-freeness is what makes the two shears differ inside the
$e_x$-directed class. It does not give cancellation on general divergence-free triples, as (e)
shows, so the tethered bracket is not Poisson even at the divergence-free shear. The witnesses use
$a=b=c$, where the Jacobiator is $3\{\{a,a\},a\}$. That is nonzero only because $T$ is symmetric
(metriplectic). The link from these Fourier coefficients to the \texttt{fderiv}-based fields of
\texttt{Torus.lean} is not formalized for the bracket.

\section{Result 2: the tethered system}\label{sec:model}
\subsection{Insertion through the projection}
The inserted system is
\begin{equation}\label{eq:tv}
\partial_t\omega+(u\cdot\nabla)\omega=(\omega\cdot\nabla)u+\nu\Delta\omega-\kappa P(|\omega|^2\omega).
\end{equation}
At velocity level it is the forced NS equation with force $F=\mathrm{curl}^{-1}(-\kappa P(|\omega|^2\omega))$.
\textbf{This is a modified Navier--Stokes system, not plain Navier--Stokes.}
The quartic weight is built in, not ansatzed:
$\T[\kappa,\|\omega\|^2,u](\omega,\omega)=-\kappa\|\omega\|^4$ for $\omega\perp u$ (\lean{quarticTether_eq}).

\begin{theorem}[Curl of forced NS]
If $(u,p)$ is a classical smooth solution of $\partial_tu+(u\cdot\nabla)u=-\nabla p+\nu\Delta u+F$ with
$\mathrm{div}\,u=0$, then $\omega=\mathrm{curl}\,u$ satisfies
$\partial_t\omega+(u\cdot\nabla)\omega=(\omega\cdot\nabla)u+\nu\Delta\omega+\mathrm{curl}\,F$.
\end{theorem}
\begin{proof}
Take the curl. Mixed partials commute on smooth fields, so $\mathrm{curl}\,\nabla p=0$ and
$\mathrm{curl}\,\Delta=\Delta\,\mathrm{curl}$, and $\partial_t$ commutes with $\mathrm{curl}$. The identity
$\mathrm{curl}((u\cdot\nabla)u)=(u\cdot\nabla)\omega-(\omega\cdot\nabla)u$ for $\mathrm{div}\,u=0$ is checked componentwise.
\end{proof}
\LeanRef{curlOfForcedNS, curl_grad, curl_laplacian, curl_advection, hasDerivAt_pd_time}

\begin{theorem}[Biot--Savart on trigonometric polynomials]
For each nonzero wavevector $k$ and amplitude $a\perp k$, the mode $b=\frac{ik\times a}{|k|^2}e^{ik\cdot x}$
satisfies $\mathrm{curl}\,b=a\,e^{ik\cdot x}$ and $\mathrm{div}\,b=0$, and it is the unique such divergence-free
mode. Summing gives a right inverse of curl on mean-zero divergence-free trigonometric polynomials.
\end{theorem}
\begin{proof}
$\mathrm{curl}\,b=\frac{ik\times(ik\times a)}{|k|^2}e^{ik\cdot x}=\frac{-(k\cdot a)k+|k|^2a}{|k|^2}e^{ik\cdot x}=a\,e^{ik\cdot x}$, and $\mathrm{div}\,b\propto k\cdot(k\times a)=0$. Curl is linear over finite sums.
\end{proof}
\LeanRef{curlC_bsC, divC_bsC, bsC_unique, curl_bsMode, div_bsMode, curl_bsSum, curlRightInverse_trig}

\begin{corollary}[Velocity form $\Rightarrow$ vorticity form]
If $\mathrm{curl}^{-1}$ is a right inverse of curl on the admissible class containing
$-\kappa P(|\omega|^2\omega)$, then every velocity solution of the tethered NS solves \eqref{eq:tv}.
\end{corollary}
\LeanRef{tetheredVorticity_of_velocity, curl_lerayTetherForce}

\subsection{$L^2$ (enstrophy) bound}
\begin{lemma}
If $P$ is a symmetric idempotent and $P\omega=\omega$, then $\ip{\omega}{-\kappa Pg}=-\kappa\ip{\omega}{g}$.
With $g=|\omega|^2\omega$ this equals $-\kappa\|\omega\|_4^4$.
\end{lemma}
\LeanRef{inner_leray_of_divFree, tether_L2_term, inner_cubic}
\begin{theorem}\label{thm:L2}
Let $E=\tfrac12\|\omega\|_2^2$ and $Q=\|\omega\|_4^4$. Assume $E'=S-\nu D-\kappa Q$ (Cited Prop.~\ref{cp:EI}),
$S\le C\sqrt{2E}\sqrt Q$ (Cited Prop.~\ref{cp:CZ}), and $(2E)^2\le Q$ (Cited Prop.~\ref{cp:Holder}). Then for $\kappa>0$,
$E(t)\le\max(E_0,\,C^2/(2\kappa^2))$ on $[0,T]$.
\end{theorem}
\begin{proof}
Young with $\lambda=\kappa/2$: $S\le\frac\kappa2 Q+\frac{C^2\cdot 2E}{2\kappa}$. Dropping $-\nu D\le0$ and using $Q\ge(2E)^2$ gives the logistic inequality
$E'\le\frac{C^2}{\kappa}E-2\kappa E^2$ (\lean{enstrophy_deriv_logistic}), so $E'<0$ whenever $E>C^2/(2\kappa^2)$.
A barrier argument (\lean{image_le_of_deriv_right_lt_deriv_boundary'}) with a level $M'$ slightly above the maximum,
followed by letting $M'\downarrow\max$, gives the bound.
\end{proof}
\LeanRef{young_lam, stretch_absorb, enstrophy_deriv_le, enstrophy_deriv_logistic, tethered_enstrophy_uniform, tethered_L2_vorticity_bounded}

\subsection{$H^1$ (palinstrophy) estimate}
\begin{lemma}[Tether is sign-definite at $H^1$]
For $\omega\in\mathbb R^3$ and $D=\nabla\omega$,
$\sum_{ij}D_{ij}\,\partial_j(|\omega|^2\omega_i)=|\omega|^2|D|^2+\tfrac12\sum_j(2\sum_k\omega_kD_{kj})^2\ge0$,
with no alignment assumption.
\end{lemma}
\LeanRef{tether_H1_density_eq, tether_H1_density_nonneg, tether_H1_term}
\begin{theorem}\label{thm:H1}
Let $G=\|\nabla\omega\|_2^2$, $L=\|\Delta\omega\|_2$. Assume $G'=2(-\nu L^2-\kappa J+N)$ with $J\ge0$
(Cited Prop.~\ref{cp:H1I}) and $N\le A\,G^{1/4}L^{3/2}$ (Cited Prop.~\ref{cp:NL}). Then
$G(t)\le G(0)\exp\!\big(\tfrac{27A^4}{16\nu^3}t\big)$.
\end{theorem}
\begin{proof}
Drop $-\kappa J\le0$. Young with exponents $(4,4/3)$: $A G^{1/4}L^{3/2}\le\nu L^2+\frac{27A^4}{32\nu^3}G$
(\lean{young_43}, \lean{amgm_four}). So $G'\le\frac{27A^4}{16\nu^3}G$, and Gr\"onwall
(\lean{le_gronwallBound_of_liminf_deriv_right_le}) finishes.
\end{proof}
\LeanRef{palin_deriv_le, tethered_palinstrophy_gronwall, higherNormControl_tethered}

\subsection{Log-loss comparison and global smoothness}
At a spatial maximum of $|\omega|$, with $y=\|\omega\|_\infty$, the maximum principle gives a nonpositive
viscous term (proved: \lean{viscousAtMax_of_max}, \lean{periodic_exists_max}). The tether gives $-\kappa y^3$
(proved from $\T$: \lean{tetherFeedback_of_quartic}). With Cited Props.~\ref{cp:SbS}, \ref{cp:BKMlog},
\ref{cp:cross} this gives
$y'\le C(1+Z_0)\,y+C\log(e+Z)\,y^2-(\kappa/2-C_{abs})y^3$ (\lean{tether_pointwise_logloss}).
The cubic dominates, so $y$ stays below the largest root of the comparison field (\lean{comparison_bound_logloss}).
\begin{theorem}[Global smoothness of the tethered system, conditional]
Assume the tethered log inequality (Cited Prop.~\ref{cp:NSlog}), higher-norm control, BKM continuation
(Cited Prop.~\ref{cp:BKM}) and local continuation (Cited Prop.~\ref{cp:LC}). Then the solution is smooth on $[0,T]$ for every $T$.
Higher-norm control is itself discharged by Theorems~\ref{thm:L2} and \ref{thm:H1} together with Cited Prop.~\ref{cp:boot}.
\end{theorem}
\begin{proof}
Let $\sigma=\sup\{T:\text{smooth on }[0,T]\}$, assuming this set is bounded. On $[0,\sigma)$ the comparison bound gives a uniform
$\sup y$. BKM continuation then gives smoothness on $[0,\sigma]$, and local continuation extends it past $\sigma$, a contradiction.
\end{proof}
\LeanRef{global_regularity_logloss, global_regularity, lemma_3_1_uniform_bound_and_continuation, frohmanian_symplectic_tether_theorem_global}

\subsection{Remaining cited propositions (explicit hypotheses, never axioms)}
Each item below is a \texttt{Prop} hypothesis in Lean. No \texttt{axiom} or \texttt{sorry} is used.
\begin{hypothesis}[\lean{TetheredEnstrophyIdentity}]\label{cp:EI} Integrated enstrophy identity for \eqref{eq:tv}. It follows from multiplying by $\omega$, integrating by parts on $T^3$, and the lemma above.\end{hypothesis}
\begin{hypothesis}[\lean{StrainL2Bound}]\label{cp:CZ} $\int\omega\cdot S\omega\le C\|\omega\|_2\|\omega\|_4^2$: Calder\'on--Zygmund / Biot--Savart, classical.\end{hypothesis}
\begin{hypothesis}[\lean{L4L2Holder}]\label{cp:Holder} $\|\omega\|_2^4\le\|\omega\|_4^4$ on the unit-volume torus (H\"older).\end{hypothesis}
\begin{hypothesis}[\lean{TetheredH1Identity}]\label{cp:H1I} The $H^1$ energy identity (Majda--Bertozzi~\cite{MB} \S3.2 for the NS part).\end{hypothesis}
\begin{hypothesis}[\lean{NonlinH1Bound}]\label{cp:NL} Sobolev/Gagliardo--Nirenberg bound of the nonlinear term~\cite{Temam,CF}.\end{hypothesis}
\begin{hypothesis}[\lean{H1ToHsBootstrap}]\label{cp:boot} $H^1$ bounds on $[0,T)$ imply $H^s$ bounds for all $s$. For plain NS this is classical parabolic regularity~\cite{Temam,CF,MB}. \textbf{The tether-specific step (controlling $\kappa P(|\omega|^2\omega)$ in $H^s$, which is plausible since $H^2\subset L^\infty$ in 3D) is not in the literature and remains to be proved.}\end{hypothesis}
\begin{hypothesis}[\lean{EnergyIdentity}, \lean{StretchingByStrain}]\label{cp:SbS} Pointwise evolution at the maximum, and stretching $\le\|S\|_\infty y^2$.\end{hypothesis}
\begin{hypothesis}[\lean{BKMLogInequality}]\label{cp:BKMlog} $\|S\|_\infty\le C(1+\|\omega\|_2+\|\omega\|_\infty\log(e+\|\omega\|_{H^s}))$~\cite{BKM}; see also Brezis--Gallou\"et and Kozono--Taniuchi~\cite{KT}.\end{hypothesis}
\begin{hypothesis}[\lean{CrossControlled}]\label{cp:cross} The cross term at the maximum is absorbed with constant $C_{abs}$.\end{hypothesis}
\begin{hypothesis}[\lean{NSTetherLogInequality}, \lean{NSTetherEnstrophyInequality}]\label{cp:NSlog} The integrated (time-continuous) form of the pointwise inequality along the solution.\end{hypothesis}
\begin{hypothesis}[\lean{BKMContinuation}]\label{cp:BKM} Beale--Kato--Majda~\cite{BKM}: bounded $\|\omega\|_\infty$ on $[0,T)$ implies continuation. Stated for plain NS; its extension to \eqref{eq:tv} is assumed.\end{hypothesis}
\begin{hypothesis}[\lean{LocalContinuation}]\label{cp:LC} Local existence for smooth periodic data (Fujita--Kato~\cite{FK}, Kato~\cite{Kato}).\end{hypothesis}
\begin{hypothesis}[\lean{CurlRightInverse} beyond trigonometric polynomials]\label{cp:BS} Proved on trigonometric polynomials. The extension to the smooth admissible class requires a density/continuity argument not yet formalized.\end{hypothesis}

\paragraph{Search of the author's own prior work for $H^1\to H^2$.}
The author's earlier repositories were searched. \texttt{NS\_Millennium\_Proof} (\texttt{Blueprint.md}, June~19, 2026, and
\texttt{Modules/AnalyticEstimates.lean}, Aug.~26, 2026), \texttt{fefferman-clay-navier-stokes}
(\texttt{HighSetDepletionAndBKM.lean}, Oct.~25, 2026) and \texttt{Clay-} (\texttt{latex/section10\_gaps.tex}, Oct.~7, 2026)
all invoke ``BKM + parabolic regularity bootstrap'' as a \emph{cited classical black box}
(in one case a \texttt{sorry}). None contains an independent proof of the $H^1\to H^2$ step, let alone the
tether-specific one. These are cited as the provenance of the plan, not as a proof.

\section{Open question, stated plainly}\label{sec:open}
\begin{enumerate}
\item \textbf{Plain Navier--Stokes is unchanged.} For a velocity field that solves plain NS, nothing above adds a
damping term. The author's October notes say so in these words:
``Because $B(F,H)=0$ by Proposition~[prop:degen], the tether does not add a damping term to the PDE (eq:vort) itself''
(\texttt{NS\_Millennium\_Proof/main.tex}, section ``Scaling Analysis and Open Analytic Questions);
``A degeneracy $B(F,H)\equiv 0$ on a kinetic-energy Hamiltonian leaves the energy field classical and does not insert
$-\kappa M^2$ into the strain'' and ``Equation (eq:missing) is not an identity of the unmodified vorticity equation''
(\texttt{certs/CubicSourceGap.tex}, 7 Oct.\ 2026);
``Derivation of the rate from the unmodified vorticity equation remains open'' (\texttt{STATUS-2026-10-10.md}).
(Bracketed labels inside quotations are the cited files' own equation references.) The mechanism is therefore the \emph{degeneracy} $\Pi_uu=0$, not a
statement that the Leray projection annihilates a velocity-level cubic source; those notes make no such claim. Here the
cubic term is inserted by hand at vorticity level through $P$, which is why the results concern
\eqref{eq:tv} only. The Lean predicates \lean{IsTrueNS} and \lean{IsTetheredNS} coincide only at $\kappa=0$
(\lean{tetheredNS_zero_iff}), and at $\kappa=0$ no feedback is realized (\lean{not_tetherRealizesFeedback_zero}).
\item \textbf{The Clay Millennium problem remains open.} Nothing in this paper or repository claims otherwise.
\item \textbf{Do real fluids carry this term?} Answering this would need either a derivation of
$-\kappa P(|\omega|^2\omega)$ from plain NS (none is known) or experimental/numerical evidence. Result~1 says
that \emph{if} such a symmetric, invariant, energy-degenerate correction exists, it must be $\T$.
\end{enumerate}

\begin{thebibliography}{9}
\bibitem{BKM} J.~T.~Beale, T.~Kato, A.~Majda, Remarks on the breakdown of smooth solutions for the 3-D Euler equations, \emph{Comm. Math. Phys.} 94 (1984) 61--66.
\bibitem{Temam} R.~Temam, \emph{Navier--Stokes Equations and Nonlinear Functional Analysis}, 2nd ed., SIAM, 1995.
\bibitem{CF} P.~Constantin, C.~Foias, \emph{Navier--Stokes Equations}, Univ.\ Chicago Press, 1988.
\bibitem{MB} A.~J.~Majda, A.~L.~Bertozzi, \emph{Vorticity and Incompressible Flow}, Cambridge, 2002.
\bibitem{KT} H.~Kozono, Y.~Taniuchi, Limiting case of the Sobolev inequality in BMO, \emph{Comm. Math. Phys.} 214 (2000) 191--200.
\bibitem{FK} H.~Fujita, T.~Kato, On the Navier--Stokes initial value problem I, \emph{Arch. Rational Mech. Anal.} 16 (1964) 269--315.
\bibitem{Kato} T.~Kato, Strong $L^p$-solutions of the Navier--Stokes equation in $\mathbb R^m$, \emph{Math. Z.} 187 (1984) 471--480.
\end{thebibliography}
\end{document}
